\section*{Abstract}
\addcontentsline{toc}{section}{Abstract}
\vspace{2cm}

Artificial intelligence has achieved impressive results in traditional board games such as Chess and Go, but modern board games have received less attention. These games can introduce additional challenges, such as random events and a wider range of possible moves, making them useful environments for studying AI behaviour in recreational game-playing contexts. This project examines how different AI approaches perform as they move from a small traditional game to the modern board game Azul, aiming to better understand their behaviours, strengths and limitations.\\

This project combines background research with practical experimentation. Existing literature on game-playing AI informed the development of three Python agents: a search-based approach using Minimax, a simulation-based approach using Monte Carlo Tree Search (MCTS), and a learning-based approach using Tabular Q-Learning (TabQL). Each approach was first implemented for tic-tac-toe and then adapted to the complete two-player game of Azul. The agents were compared through head-to-head simulated matches.\\

In tic-tac-toe matches, every game ended in a draw. In Azul, MCTS won more games than Minimax, while both won every reported match against TabQL. These results suggest that the approaches performed differently under the configurations tested, but the limited number of matches prevents broader conclusions about playing strength. TabQL’s weaker Azul results may reflect the difficulty of learning in a game with many possible states and moves.\\

This research provides a foundation for applying AI techniques to a wider range of modern board games. Future work could investigate other reinforcement-learning methods, explore hybrid approaches, human-interaction studies, and applications in game design and decision-making contexts.